\documentclass[11pt]{article}

\usepackage[T1]{fontenc}
\usepackage[utf8]{inputenc}
\usepackage{lmodern}
\usepackage{microtype}
\usepackage[a4paper,margin=25mm]{geometry}
\usepackage{amsmath,amssymb,mathtools,bm}
\usepackage{booktabs,longtable,tabularx,array}
\usepackage{enumitem}
\usepackage{xcolor}
\usepackage{hyperref}

\hypersetup{
  colorlinks=true,
  linkcolor=blue!55!black,
  citecolor=blue!55!black,
  urlcolor=blue!55!black,
  pdftitle={Theory and implementation of reusable real-time TDDFT in Aion},
  pdfauthor={Aion project}
}

\newcommand{\ii}{\mathrm{i}}
\newcommand{\dd}{\mathrm{d}}
\newcommand{\ee}{\mathrm{e}}
\newcommand{\hH}{H_{\mathrm{eom}}}
\newcommand{\Pzero}{\mathrm{P0}}
\newcommand{\Eone}{\mathrm{E1}}
\newcommand{\ReTr}{\operatorname{Re}\operatorname{Tr}}
\newcommand{\ImTr}{\operatorname{Im}\operatorname{Tr}}
\newcommand{\Herm}{\operatorname{Herm}}
\newcommand{\Tr}{\operatorname{Tr}}
\newcommand{\norm}[1]{\left\lVert#1\right\rVert}
\newcommand{\vect}[1]{\bm{#1}}
\newcommand{\code}[1]{\texttt{#1}}

\newenvironment{statusbox}[1]
  {\begin{quote}\small\textbf{Status --- #1.}\ }
  {\end{quote}}

\title{Theory and implementation of reusable real-time TDDFT in Aion\\
\large Bare length and velocity gauges, covariant P0/P0+E1 dynamics,\\
currents, energies, and connection-aware propagation}
\author{Aion project}
\date{8 September 2026}

\begin{document}

\maketitle

\begin{abstract}
This note fixes the physical definitions, sign and orientation conventions,
equations of motion, observables, numerical algorithms, and software boundaries
for the proposed Aion 0.2 real-time time-dependent density-functional theory
(RT-TDDFT) implementation. It treats four formulations on one common
finite-basis electronic reference: bare length gauge, bare velocity gauge,
covariant P0, and covariant P0+E1. Particular attention is paid to quantities
that cannot be transferred uncritically between formulations: the primary
current, the mechanical-current diagnostic, matter and generator energies,
source work, and the evolving-metric connection. The production propagation
algorithm is self-consistent exponential midpoint (SCEM), with a
connection-aware transport for the covariant formulations and no
time-dependent L\"owdin frame. The final sections map these definitions to a
reusable API, storage schema, and validation hierarchy. This is primarily a
normative design document. Its non-numerical package, configuration, identity,
schema, status, and API contracts were implemented in WP1; the physical and
numerical layers remain gated future work unless explicitly identified.
\end{abstract}

\tableofcontents

\section{Document status and scope}

\subsection{Meaning of status labels}

Four kinds of statements are distinguished throughout this note.

\begin{description}[leftmargin=34mm,style=nextline]
\item[Established theory]
Standard RT-TDDFT, minimal coupling, gauge transformation, evolving-Hilbert-
space geometry, and rational approximation results supported by the cited
literature.
\item[Normative Aion 0.2]
An agreed definition or software/numerical contract that the refactor must
implement.
\item[Draft-validated]
Physics or algorithms exercised in the present research tree, but not yet
available through the clean 0.2 API or fully covered by its tests.
\item[Planned/unvalidated]
Agreed work whose complete reusable implementation and validation remain to
be done.
\end{description}

\begin{statusbox}{Normative Aion 0.2}
When this note uses \emph{must} or gives an unqualified stored name, it
specifies the intended 0.2 contract. The implementation plan in
\path{docs/refactor_implementation_plan.md} controls the order and acceptance
gates. That plan is deliberately mutable; changes to the physics conventions
require an explicit reviewed decision.
\end{statusbox}

\subsection{Validated physical domain}

The first release targets finite, fixed-nucleus, all-electron molecules in
restricted spin-summed Kohn--Sham DFT, with pure LDA or GGA functionals and
prescribed spatially uniform electric fields. Atomic units are used internally
and the numerical precision is double real/complex precision on CPU and GPU.
The nuclear attraction is local.

The following are explicitly outside the first physics stack:

\begin{itemize}
\item pseudopotentials, effective-core potentials, and nonlocal ionic
  operators;
\item hybrid and exact-exchange current corrections;
\item unrestricted, noncollinear, or relativistic spin dynamics;
\item moving nuclei and their basis-velocity connection;
\item periodic boundary conditions;
\item Maxwell backreaction and adaptive time grids.
\end{itemize}

These cases must be rejected rather than silently treated with formulas valid
only for local all-electron Hamiltonians. The source and observable interfaces
are nevertheless designed so that spatially varying fields, aloof relativistic
projectiles, EELS, and possibly periodic velocity gauge can be added later
without redefining the present uniform-field theory. They are extension seams,
not present work packages.

\section{Notation and conventions}

\subsection{Atomic basis and density}

Let \(\{\lvert\phi_\mu(t)\rangle\}\) be a generally time-dependent
nonorthogonal atomic-orbital (AO) basis. The overlap and occupied orbitals are
\begin{equation}
 S_{\mu\nu}(t)=\langle\phi_\mu(t)\vert\phi_\nu(t)\rangle,
 \qquad
 \lvert\psi_i(t)\rangle
 =\sum_\mu\lvert\phi_\mu(t)\rangle C_{\mu i}(t).
\end{equation}
With occupations collected in \(f\), the canonical AO density convention is
\begin{equation}
 P=CfC^\dagger,\qquad
 N=\ReTr(PS),\qquad
 \langle O\rangle=\ReTr(PO).
 \label{eq:density-convention}
\end{equation}
The occupations contain the complete spin multiplicity. No contraction may
insert an implicit closed-shell factor of two. A different covariant/dual or
orthonormal representation is a different type and requires an explicit
conversion.

Electronic charge and mass are kept symbolically as \(q\) and \(m\);
electrons have \(q=-1\), \(m=1\) in atomic units. We retain \(\hbar\)
in derivations to expose phases and signs, although \(\hbar=1\) in the
implementation. The reduced vector potential is
\begin{equation}
 \vect a=\frac{\vect A}{c},
\end{equation}
and the full stored name is \code{vector\_potential\_reduced}. This avoids
an otherwise persistent ambiguity about factors of \(c\).

\subsection{Coordinates, origin, and molecular observables}

Input coordinates are preserved exactly: Aion must never silently recenter a
molecule. An explicit electromagnetic origin \(\vect O\), defaulting to the
input-coordinate origin, is stored. Electronic, fixed-nuclear, and total
dipoles use this same origin:
\begin{align}
 \vect\mu_{\mathrm e}
   &=q\,\ReTr\!\left[P(\vect r-\vect O S)\right],\\
 \vect\mu_{\mathrm N}
   &=\sum_A Z_A(\vect R_A-\vect O),\\
 \vect\mu_{\mathrm{tot}}
   &=\vect\mu_{\mathrm e}+\vect\mu_{\mathrm N}.
\end{align}
For the finite molecule, \(\vect J\) denotes total charge current, not a
volume-normalized current density.

\subsection{Time grid and sampling location}

For a fixed step, integer index is authoritative:
\begin{equation}
 t_n=t_0+n\Delta t,\qquad n=0,\ldots,N.
\end{equation}
There are \(N\) propagation intervals and \(N+1\) endpoint states.
Floating-point time is stored as a convenience. State, dipole, current, and
sampled energy records are endpoint-centred. Work increments, point power,
and Ward residuals associated with a propagation interval are midpoint- or
interval-centred and carry their own step/time arrays.

\section{Electronic structure and the AO equation of motion}

\subsection{Field-free Kohn--Sham model}

The all-electron molecular reference is a converged restricted Kohn--Sham
calculation. For a pure LDA/GGA functional its field-free Fock/Kohn--Sham
matrix is written
\begin{equation}
 H_0[P]=T+V_{\mathrm{nuc}}+J[P]+V_{\mathrm{xc}}[P].
\end{equation}
The corresponding field-free energy decomposition is
\begin{align}
 T_s[P]&=\ReTr(PT),&
 E_{\mathrm{eN}}[P]&=\ReTr(PV_{\mathrm{nuc}}),\\
 E_{\mathrm H}[P]&=\frac12\ReTr(PJ[P]),&
 E_{\mathrm{xc}}[P]&=E_{\mathrm{xc}}[n_P],\\
 E_{\mathrm{NN}}&=\sum_{A<B}\frac{Z_AZ_B}{R_{AB}},&
 U_0[P]&=T_s+E_{\mathrm{eN}}+E_{\mathrm H}
       +E_{\mathrm{xc}}+E_{\mathrm{NN}}.
 \label{eq:dft-energy}
\end{align}
Real-time propagation is used within the Runge--Gross density--potential
framework\cite{RungeGross1984}; PBE is a representative pure GGA in the
validated functional class\cite{PBE1996}.
The implementation uses PySCF for the RKS reference, numerical grid, effective
potential, and functional energy\cite{Sun2018}. The first release stores
combined exchange--correlation energy; an exchange/correlation split is not a
required observable.

\subsection{Connection in a time-dependent nonorthogonal basis}

Projection of the time-dependent Kohn--Sham equation into an evolving basis
gives
\begin{equation}
 \ii\hbar\left(S\dot C+\omega C\right)=\hH C,
 \qquad
 \omega_{\mu\nu}
 =\langle\phi_\mu\vert\dot\phi_\nu\rangle .
 \label{eq:general-eom}
\end{equation}
The geometry is not an optional correction. It imposes metric compatibility
\begin{equation}
 \dot S=\omega+\omega^\dagger,
 \label{eq:metric-compatibility}
\end{equation}
which guarantees preservation of the physical overlap norm under the exact
equation. Artacho and O'Regan discuss why an evolving Hilbert space is
geometrically different from a standard fixed-coordinate matrix
ODE\cite{ArtachoORegan2017}.

For implementation, define
\begin{equation}
 \mathcal A=S^{-1}\left[-\omega-\frac{\ii}{\hbar}\hH\right],
 \qquad \dot C=\mathcal A C,
 \qquad \dot P=\mathcal A P+P\mathcal A^\dagger.
 \label{eq:generator-density}
\end{equation}
Knowledge of \(S(t)\) and \(\omega(t)\) at arbitrary times removes any need
to integrate a separate metric equation. It does not remove the need to
transport the electronic state consistently between two different metrics,
nor the need to solve the nonlinear midpoint Hamiltonian self-consistently.

\begin{statusbox}{Normative Aion 0.2}
Every formulation returns the explicit triple
\((S,\hH,\omega)\). The older implicit split between an intrinsic
Hamiltonian and a residual connection is not a public contract.
\end{statusbox}

\section{Electromagnetic sources and gauge data}

\subsection{Potential-first representation}

The minimal-coupling one-particle operator is
\begin{equation}
 h(t)=\frac{1}{2m}\left(\vect p-q\vect a(\vect r,t)\right)^2
      +q\phi(\vect r,t)+v_0,
 \label{eq:minimal-coupling}
\end{equation}
where \(\vect p=-\ii\hbar\nabla\). The electric field is
\begin{equation}
 \vect E=-\partial_t\vect a-\nabla\phi.
\end{equation}
A gauge function \(\chi\) acts as
\begin{equation}
 \vect a' = \vect a+\nabla\chi,\qquad
 \phi'=\phi-\dot\chi,\qquad
 \psi'=\exp(\ii q\chi/\hbar)\psi.
 \label{eq:gauge-transform}
\end{equation}

The reusable source layer describes the physical field separately from a
chosen gauge representation. It is potential-first and evaluable at arbitrary
times, including every SCEM midpoint. A field-first convenience input is
compiled to consistent potentials before propagation; \(\vect a\) is never
constructed by incremental numerical integration inside the electronic hot
loop.

\subsection{Projected node/link sample}

For the covariant formulation, associate AO \(\mu\) with an atomic anchor
\(a(\mu)\). The compact projected source sample is
\begin{equation}
 \left\{\phi_a(t),L_{ab}(t),\dot L_{ab}(t)\right\},\qquad
 L_{ab}=\int_{\vect R_b}^{\vect R_a}\vect a\cdot\dd\vect\ell,
 \quad a<b .
 \label{eq:source-sample}
\end{equation}
The pair electromotive potential is
\begin{equation}
 V_{ab}=-\dot L_{ab}-\phi_a+\phi_b.
 \label{eq:pair-emf}
\end{equation}
Links and EMFs are oriented \(b\to a\), while \(I_{ab}>0\) denotes
charge flow \(a\to b\). With the corresponding signed incidence matrix,
\begin{align}
 \dot Q_a+\sum_b I_{ab}&=0,\label{eq:continuity}\\
 \vect J_{\mathrm{graph}}
   &=\sum_{a<b}I_{ab}(\vect R_b-\vect R_a),\label{eq:graph-current}\\
 \mathcal P_{\mathrm{source}}
   &=-\sum_{a<b}I_{ab}V_{ab}.\label{eq:graph-power}
\end{align}
Under Eq.~\eqref{eq:gauge-transform},
\(L'_{ab}=L_{ab}+\chi_a-\chi_b\), so Eq.~\eqref{eq:pair-emf} is invariant.

Multiple physical providers are added at the primitive potential/link level
before projection or Wilson dressing. A temporal gate \(g(t)\) acts on
potentials,
\begin{equation}
 \phi^g=g\phi,\qquad L^g=gL,\qquad
 \dot L^g=\dot g L+g\dot L.
\end{equation}
Omitting the product-rule term would omit a real switching field.

\subsection{Finite pulses and discrete kicks}

A compactly supported resonant pulse may be defined in the vector potential:
\begin{equation}
 \vect a(t)=a_0
 \sin^2\!\left(\frac{\pi t}{T}\right)
 \sin(\Omega t+\varphi)\,\hat{\vect e},
 \qquad 0\le t\le T,
 \label{eq:finite-pulse}
\end{equation}
with \(T\) an exact integer number of carrier cycles and
\(\vect E=-\dot{\vect a}\). The amplitude is normalized to the requested
peak electric field. Equation~\eqref{eq:finite-pulse} has
\(\vect a(0)=\vect a(T)=0\) and zero net impulse. This is a property of this
pulse family, not a restriction on the physical use of pulses.

If a maximum step is requested, choose
\begin{equation}
 N=\left\lceil T/\Delta t_{\max}\right\rceil,\qquad
 \Delta t=T/N
\end{equation}
so the boundary is exactly a grid point. General pulses may carry nonzero
impulse. A simulation or analysis workflow may request a zero-impulse
constraint when its scientific purpose requires one; the pulse abstraction
itself imposes no such policy.

A delta kick is a uniquely identified discrete event, not a narrow sampled
pulse. Events lie exactly on grid boundaries and are stored with pre- and
post-event records at the same physical time, distinguished by side and
sequence. The exact map is formulation-owned:

\begin{itemize}
\item bare LG applies the coefficient-space representation of
  \(\exp(\ii\vect\kappa\cdot\vect r)\);
\item bare VG leaves coefficients continuous and makes the discontinuous
  change \(\vect a(0^+)=-\vect\kappa\) in the adopted electronic sign
  convention;
\item P0+E1 applies its covariant gauge/state event map.
\end{itemize}

Applied event identifiers are checkpointed, making a restarted event
idempotent.

\section{Bare length-gauge formulation}

\subsection{Equation of motion}

For a uniform field in length gauge, \(\vect a=0\) and
\(\phi=-\vect E\cdot(\vect r-\vect O)\). In a fixed AO basis,
\begin{equation}
 S=S_0,\qquad \omega=0,\qquad
 H_{\mathrm{LG}}[P,t]
 =H_0[P]-q\vect E(t)\cdot(\vect r-\vect O S).
 \label{eq:lg-hamiltonian}
\end{equation}

\subsection{Dipole and currents}

The primary current is the analytic time derivative of the length-gauge
dipole:
\begin{align}
 \mu_{\mathrm e,\alpha}
   &=q\ReTr[P R_\alpha],
 &R_\alpha&=r_\alpha-O_\alpha S,\\
 J_{\mathrm{LG},\alpha}
   &=q\ReTr[\dot P R_\alpha]\nonumber\\
   &=\Re\left\{
     \frac{\ii q}{\hbar}
     \Tr\left[P\left(
       H_{\mathrm{LG}}S^{-1}R_\alpha
       -R_\alpha S^{-1}H_{\mathrm{LG}}
     \right)\right]\right\}.
 \label{eq:lg-current}
\end{align}
The second line follows from Eq.~\eqref{eq:generator-density} and is not a
finite difference.

For comparison, a local all-electron ambient mechanical diagnostic is
\begin{equation}
 J_{\mathrm{para},\alpha}
 =\frac q m\ReTr(Pp_\alpha),\qquad
 J_{\mathrm{dia},\alpha}=0,\qquad
 \vect J_{\mathrm{mech}}=\vect J_{\mathrm{para}}.
 \label{eq:lg-mechanical}
\end{equation}
The analytic dipole current is the primary LG observable. In a finite basis,
Eq.~\eqref{eq:lg-current} and Eq.~\eqref{eq:lg-mechanical} need not coincide;
their difference measures finite-subspace failure of continuum commutator
identities.

\subsection{Energy}

The length-gauge matter energy excludes the source coupling:
\begin{equation}
 U_{\mathrm{matter}}^{\mathrm{LG}}=U_0[P].
\end{equation}
The electronic and fixed-nuclear scalar couplings and generator energy are
\begin{align}
 U_{\mathrm{EM,e}}^{\mathrm{LG}}
   &=-\vect E\cdot\vect\mu_{\mathrm e},&
 U_{\mathrm{EM,N}}^{\mathrm{LG}}
   &=-\vect E\cdot\vect\mu_{\mathrm N},\\
 E_{\mathrm{generator}}^{\mathrm{LG}}
   &=U_{\mathrm{matter}}^{\mathrm{LG}}
     +U_{\mathrm{EM,e}}^{\mathrm{LG}}
     +U_{\mathrm{EM,N}}^{\mathrm{LG}}.
\end{align}
For fixed nuclei,
\begin{align}
 \frac{\dd U_{\mathrm{matter}}^{\mathrm{LG}}}{\dd t}
   &=\vect E\cdot\vect J_{\mathrm{LG}},\label{eq:lg-work}\\
 \frac{\dd E_{\mathrm{generator}}^{\mathrm{LG}}}{\dd t}
   &=-\dot{\vect E}\cdot\vect\mu_{\mathrm{tot}}.
 \label{eq:lg-generator-rate}
\end{align}
These are different valid energies. The phrase \emph{total energy} is
therefore insufficient without the qualifier \code{matter} or
\code{generator}. When the field vanishes at both selected endpoints, the
final deposited matter energy and final generator-energy gain agree.

\section{Bare velocity-gauge formulation}

\subsection{Hamiltonian}

For \(\phi=0\) and spatially uniform \(\vect a(t)\), expansion of
Eq.~\eqref{eq:minimal-coupling} gives the AO matrix
\begin{equation}
 H_{\mathrm{VG}}[P,t]
 =H_0[P]
 -\frac q m\vect a(t)\cdot\vect p
 +\frac{q^2}{2m}|\vect a(t)|^2S.
 \label{eq:vg-hamiltonian}
\end{equation}
This is the local all-electron velocity-gauge form used in the NAO
RT-TDDFT treatment of Pemmaraju \emph{et al.}\cite{Pemmaraju2018}. For a
nonlocal ionic potential, an additional gauge-transformed nonlocal operator
and corresponding current correction are required; they are outside scope.

With PySCF derivative overlaps, the implementation fixes and tests the
canonical-momentum convention
\begin{equation}
 p_\alpha=\ii\hbar
 \langle\nabla_\alpha\phi_\mu\vert\phi_\nu\rangle,
\end{equation}
including its Hermiticity and sign.

\subsection{Mechanical current}

The bare VG primary current is the mechanical current,
\begin{align}
 J_{\mathrm{para},\alpha}
   &=\frac q m\ReTr(Pp_\alpha),\label{eq:vg-jpara}\\
 J_{\mathrm{dia},\alpha}
   &=-\frac{q^2}{m}N a_\alpha,\label{eq:vg-jdia}\\
 J_{\mathrm{mech},\alpha}
   &=J_{\mathrm{para},\alpha}+J_{\mathrm{dia},\alpha}.
 \label{eq:vg-current}
\end{align}
The two pieces must be available separately. Plotting only
\(q\ReTr(Pp_\alpha)\) as the VG physical current omits the diamagnetic
piece.

\subsection{Mechanical kinetic and matter energy}

The mechanical kinetic energy is
\begin{align}
 T_{\mathrm{mech}}[P,t]
 &=T_{\mathrm{can}}[P]+T_{\mathrm{linear}}[P,t]
   +T_{\mathrm{dia}}[P,t],\label{eq:vg-tkin}\\
 T_{\mathrm{can}}&=\ReTr(PT),\nonumber\\
 T_{\mathrm{linear}}
   &=-\frac q m\vect a\cdot\ReTr(P\vect p),\nonumber\\
 T_{\mathrm{dia}}
   &=\frac{q^2}{2m}N|\vect a|^2.\nonumber
\end{align}
The VG matter energy is
\begin{equation}
 U_{\mathrm{matter}}^{\mathrm{VG}}
 =T_{\mathrm{mech}}+E_{\mathrm{eN}}+E_{\mathrm H}
  +E_{\mathrm{xc}}+E_{\mathrm{NN}}.
 \label{eq:vg-energy}
\end{equation}
In the pure VG representation it is also the electronic generator energy for
the prescribed field, and
\begin{equation}
 \frac{\dd U_{\mathrm{matter}}^{\mathrm{VG}}}{\dd t}
 =\vect E\cdot\vect J_{\mathrm{mech}}.
 \label{eq:vg-work}
\end{equation}
An energy that retains only \(T_{\mathrm{can}}\) contains a large,
apparently anomalous oscillatory error during the pulse. Both the linear
\(\vect a\cdot\vect p\) and quadratic \(a^2\) contributions in
Eq.~\eqref{eq:vg-tkin} are compulsory.

\section{Covariant P0 formulation}

\subsection{Wilson-dressed overlap}

Let each AO carry an anchor position \(\vect R_{a(\mu)}\). With the link
oriented from the anchor of \(\nu\) to that of \(\mu\), define
\begin{equation}
 \Theta_{\mu\nu}(t)
 =\exp\!\left[\frac{\ii q}{\hbar}
   L_{a(\mu)a(\nu)}(t)\right],
 \qquad
 S_{\mu\nu}(t)=\Theta_{\mu\nu}(t)S^0_{\mu\nu}.
 \label{eq:p0-overlap}
\end{equation}
The inverse-dressed density passed to the field-free electronic-structure
model is
\begin{equation}
 P^0=\Theta^*\odot P,
 \label{eq:inverse-density}
\end{equation}
where \(\odot\) is elementwise multiplication. If
\(H_0[P^0]\) is constructed by PySCF, the P0 equation-of-motion matrix is
obtained by the forward dressing
\begin{equation}
 H_{\Pzero}[P,t]=\Theta\odot H_0[P^0].
 \label{eq:p0-hamiltonian}
\end{equation}

\subsection{Projected connection}

Introduce the diagonal site scalar-potential connection
\begin{equation}
 \sigma_{\mu\nu}
 =\delta_{\mu\nu}\frac{\ii q}{\hbar}
   \phi_{a(\mu)}.
\end{equation}
The site-covariant overlap derivative and P0 connection are
\begin{align}
 D_{\mathrm{site}}S
   &=\dot S+\sigma S-S\sigma,\label{eq:site-derivative}\\
 \omega_{\Pzero}
   &=S\sigma+\frac12D_{\mathrm{site}}S.
 \label{eq:p0-connection}
\end{align}
Because \(\sigma^\dagger=-\sigma\) for a real scalar potential,
Eq.~\eqref{eq:p0-connection} satisfies
\(\omega_{\Pzero}+\omega_{\Pzero}^\dagger=\dot S\) exactly.

The P0 equation is consequently
\begin{equation}
 \ii\hbar(S\dot C+\omega_{\Pzero}C)
 =H_{\Pzero}[P,t]C .
 \label{eq:p0-eom}
\end{equation}
The scalar-potential and link time dependence is geometrical: it is carried by
the projected connection and metric, not added again to
Eq.~\eqref{eq:p0-hamiltonian}.

\subsection{Site charge and graph current}

Let \(M_a\) be the fixed AO selector/projector associated with anchor
\(a\). The symmetrized source charge is
\begin{equation}
 Q_a=\frac q2\ReTr\!\left[P\{M_a,S\}\right].
 \label{eq:site-charge}
\end{equation}
For each stored pair, let \(S_{ab}\) denote its oriented overlap contribution
and \(H_{\mathrm{source},ba}\) the corresponding explicit source block.
Define
\begin{equation}
 K_{\mathrm{dyn}}
 =H_{\mathrm{dyn}}-\frac{\ii\hbar}{2}D_{\mathrm{site}}S.
\end{equation}
Here \(H_{\mathrm{dyn}}\) is the Hermitian matrix in the equivalent
coefficient-velocity equation after the non-metric connection terms have been
moved to its right-hand side. Thus \(H_{\mathrm{dyn}}=H_{\Pzero}\) for P0 and
\(H_{\mathrm{dyn}}=H_{\Pzero}+V_{\Eone}\) for P0+E1. This algebraic definition
does not alter the normative connection split in Eq.~\eqref{eq:e1-split}.
The variational graph-current expression used by the covariant model is
\begin{align}
 I_{ab}
 &=\frac q\hbar\ImTr\!\left[
 P\left(
 S_{ab}S^{-1}K_{\mathrm{dyn}}
 +K_{\mathrm{dyn}}^\dagger S^{-1}S_{ab}
 \right)\right]\nonumber\\
 &\quad
 +\frac{2q}{\hbar}\ImTr\!\left(PH_{\mathrm{source},ba}\right).
 \label{eq:p0-pair-current}
\end{align}
Together with Eqs.~\eqref{eq:continuity}--\eqref{eq:graph-power}, this gives
the P0 charge, current vector, and source power. The pair orientation and the
definition of the source block are part of the data model, not choices made by
individual analysis scripts.

\section{P0+E1 correction}

\subsection{Central dipoles and variational connection}

For the AO pair \((\mu,\nu)\), define the anchor-pair centre
\begin{equation}
 \vect R_{\mu\nu}
 =\frac12\left(\vect R_{a(\mu)}+\vect R_{a(\nu)}\right)
\end{equation}
and field-free central dipole matrices
\begin{equation}
 d^\alpha_{\mu\nu}
 =q\left[r^\alpha_{\mu\nu}
          -R^\alpha_{\mu\nu}S^0_{\mu\nu}\right].
 \label{eq:central-dipole}
\end{equation}
Their dressed Hermitian form and E1 variational potential are
\begin{align}
 D_\alpha(t)
   &=\Herm\!\left[\Theta(t)\odot d^\alpha\right],
 &\Herm X&=\frac12(X+X^\dagger),\label{eq:dressed-dipole}\\
 V_{\Eone}[P,t]
   &=-\sum_\alpha E_\alpha(t)D_\alpha(t).
 \label{eq:e1-potential}
\end{align}

The defining physics split is
\begin{equation}
 \boxed{
 \hH=H_{\Pzero},\qquad
 \omega_{\Pzero+\Eone}
 =\omega_{\Pzero}+\frac{\ii}{\hbar}V_{\Eone}.}
 \label{eq:e1-split}
\end{equation}
Putting \(V_{\Eone}\) in both \(\hH\) and \(\omega\) double counts it.
Equation~\eqref{eq:e1-split} expresses the projection of the ambient
electromagnetic connection into the P0+E1 equivariant subspace.

\subsection{Covariant dipole}

Writing \(N_a=Q_a/q\), the electronic covariant dipole is
\begin{equation}
 \mu_{\Pzero+\Eone,\alpha}
 =q\sum_a N_aR_{a\alpha}+\ReTr(PD_\alpha).
 \label{eq:covariant-dipole}
\end{equation}
The first term transports site charge between anchors; the second restores
intra-/inter-site polarization retained at E1.

\subsection{E1 residual and total source current}

The reduced-vector-potential derivative of the dressed central dipole is
\begin{equation}
 D_{\alpha;\beta}
 =\frac{\partial D_\alpha}{\partial a_\beta}
 =\Herm\!\left[
 \frac{\ii q}{\hbar}
 (R_{a(\mu),\beta}-R_{a(\nu),\beta})
 \Theta_{\mu\nu}d^\alpha_{\mu\nu}
 \right].
 \label{eq:dressed-dipole-derivative}
\end{equation}
The E1 residual source current is
\begin{equation}
 J_{\Eone,\beta}
 =\frac{\dd}{\dd t}\ReTr(PD_\beta)
 +\sum_\alpha E_\alpha\ReTr(PD_{\alpha;\beta}).
 \label{eq:e1-current}
\end{equation}
Its first term is the intrinsic polarization rate and its second term is the
phase response. Both pieces are recorded.

The total P0+E1 variational source current is
\begin{equation}
 \vect J_{\mathrm{src}}
 =\sum_{a<b}I_{ab}^{(\Pzero\mid\Eone)}
  (\vect R_b-\vect R_a)+\vect J_{\Eone}.
 \label{eq:source-current}
\end{equation}
In Eq.~\eqref{eq:p0-pair-current}, the coefficient/dynamical part sees the
full P0+E1 equation, while the explicit source-Hamiltonian block is the base
P0 source contribution. This separation prevents E1 double counting.

The defining checks are
\begin{align}
 \dot Q_a+\sum_b I_{ab}&=0,\label{eq:cov-continuity}\\
 \vect J_{\mathrm{src}}
   &=\frac{\dd\vect\mu_{\Pzero+\Eone}}{\dd t},\label{eq:j-mu}\\
 \frac{\dd U_{\mathrm{cov}}}{\dd t}
   &=\vect E\cdot\vect J_{\mathrm{src}}.
 \label{eq:cov-ward}
\end{align}

\begin{statusbox}{Draft-validated}
The current research tree contains the P0/P0+E1 tensors, connection split,
site charges, pair-current expression, and analytic E1-current pieces. Their
clean typed integration with the complete energy ledger and storage schema is
planned, not yet complete.
\end{statusbox}

\section{Mechanical-current diagnostics}

\subsection{Generic local contraction}

For an arbitrary density and prescribed local vector potential, define AO
matrices
\begin{equation}
 A^{\mathrm{op}}_{\alpha,\mu\nu}(t)
 =\langle\phi_\mu\vert
 a_\alpha(\vect r,t)\vert\phi_\nu\rangle.
\end{equation}
The reusable low-level diagnostic is
\begin{align}
 J_{\mathrm{para},\alpha}
   &=\frac q m\ReTr(Pp_\alpha),\label{eq:generic-jpara}\\
 J_{\mathrm{dia},\alpha}
   &=-\frac{q^2}{m}\ReTr(PA^{\mathrm{op}}_\alpha),
     \label{eq:generic-jdia}\\
 \vect J_{\mathrm{mech}}
   &=\vect J_{\mathrm{para}}+\vect J_{\mathrm{dia}}.
 \label{eq:generic-jmech}
\end{align}
For a uniform field,
\(A^{\mathrm{op}}_\alpha=a_\alpha S\), giving
Eq.~\eqref{eq:vg-jdia}. The density need only be finite, Hermitian, and
shape-compatible; idempotency, positivity, and electron count are optional
diagnostics rather than hidden preconditions.

\subsection{Covariant ambient projection}

In the uniform covariant representation the canonical momentum and
diamagnetic operator are
\begin{align}
 p^{\mathrm{can,cov}}_\alpha
 &=\Theta\odot(p^0_\alpha+q a_\alpha S^0),
 \label{eq:cov-pcan}\\
 p^{\mathrm{dia,cov}}_\alpha
 &=-q a_\alpha S(t).
\end{align}
Their explicit \(a_\alpha\) terms cancel, so the total ambient-projected
mechanical operator reduces to
\begin{equation}
 p^{\mathrm{mech,cov}}_\alpha
 =\Theta\odot p^0_\alpha.
 \label{eq:cov-pmech}
\end{equation}
The resulting observable is named
\code{ambient\_projected\_mechanical\_current}. It is a useful common
diagnostic, but it is not the P0+E1 variational source current in
Eq.~\eqref{eq:source-current}. The two must never share a generic stored
\code{current} label.

\section{Covariant energy and universal energy ledger}

\subsection{Covariant matter energy}

The gauge-invariant P0/P0+E1 matter energy is evaluated from the inverse-
dressed density in Eq.~\eqref{eq:inverse-density}:
\begin{equation}
 U_{\mathrm{cov}}[P,t]
 =T_s[P^0]+E_{\mathrm{eN}}[P^0]+E_{\mathrm H}[P^0]
  +E_{\mathrm{xc}}[P^0]+E_{\mathrm{NN}}.
 \label{eq:cov-energy}
\end{equation}
The E1 source-coupling expectation
\begin{equation}
 U_{\Eone,\mathrm{coupling}}=\ReTr(PV_{\Eone})
\end{equation}
is stored separately and is not included in
Eq.~\eqref{eq:cov-energy}. This is necessary for the Ward identity
Eq.~\eqref{eq:cov-ward}. Equivalent dressed-frame contractions are valuable
independent validation identities, not alternative user-facing definitions.

\subsection{Qualified totals}

The energy observable records, when meaningful for a formulation:

\begin{enumerate}
\item canonical/paramagnetic kinetic, vector-potential-linear kinetic,
  diamagnetic kinetic, and their mechanical sum;
\item electron--nuclear, Hartree, combined exchange--correlation, and
  nuclear--nuclear energies;
\item electronic scalar EM coupling, fixed-nuclear scalar EM coupling, and
  their sum;
\item \code{energy\_matter\_total};
\item \code{energy\_generator\_total};
\item absorbed energy and accumulated source work;
\item analytic energy rate, source power, and Ward residuals.
\end{enumerate}

The conventional DFT word \emph{total} is retained only with a qualifier.
There is no public \code{total\_energy}, because
Eqs.~\eqref{eq:lg-work}, \eqref{eq:lg-generator-rate},
\eqref{eq:vg-work}, and \eqref{eq:cov-ward} show that different totals
answer different physical questions during a pulse.

For every formulation,
\begin{equation}
 \Delta U_{\mathrm{abs}}(t)
 =U_{\mathrm{matter}}(t)-U_{\mathrm{matter}}(t_0)
 \label{eq:absorbed-energy}
\end{equation}
is compared to independently accumulated source work. Fixed nuclei have no
current and do not absorb power; their scalar coupling is nevertheless stored
to establish origin cancellation in a neutral system.

\subsection{Independent analytic energy derivative}

For the pure-functional DFT energy, functional differentiation gives
schematically
\begin{equation}
 \frac{\dd U}{\dd t}
 =\ReTr\!\left[F_{\mathrm{energy}}[P]\,\dot P\right]
 +\left.\frac{\partial U}{\partial t}\right|_P,
 \label{eq:analytic-energy-rate}
\end{equation}
where \(\dot P\) is obtained analytically from
Eq.~\eqref{eq:generator-density}. This route is independent of the
\(\vect E\cdot\vect J\) contraction and therefore provides a
non-tautological Ward test. Finite differences of sampled energies are only
secondary convergence diagnostics.

At every accepted propagation interval, source work is advanced from the
converged midpoint:
\begin{equation}
 W_{n+1}=W_n+\Delta t\,
 \mathcal P_{\mathrm{source}}(t_{n+1/2}) .
 \label{eq:work-quadrature}
\end{equation}
The midpoint state and source already exist, so no additional Fock build is
required. Work is accumulated every interval even when energy is not sampled
there. A robust normalized Ward residual reports both absolute and relative
scales and should decrease quadratically with \(\Delta t\) in the asymptotic
regime.

\begin{statusbox}{Planned/unvalidated}
The draft tree has partial aggregate energy methods and ad hoc reconstruction
of missing bare-VG kinetic terms. The complete qualified ledger and analytic
functional energy derivative described here are not yet a reusable
implementation.
\end{statusbox}

\section{Self-consistent exponential midpoint propagation}

\subsection{Why the connection matters numerically}

For a fixed linear ODE, exponentiating a midpoint generator is routine. Here
three facts interact:

\begin{enumerate}
\item the density determines the Kohn--Sham Hamiltonian;
\item the coordinate metric differs at the beginning and end of an interval;
\item the connection contains both an analytically integrable diagonal site
  piece and a residual projected piece that belongs in the dynamical
  generator.
\end{enumerate}

An exponential that is unitary in one Euclidean or L\"owdin-orthogonalized
frame need not map the \(S_n\) norm to the \(S_{n+1}\) norm. Small defects can
accumulate and feed back into the nonlinear density. The remedy is to respect
parallel transport and cross-metric geometry explicitly, not to pretend that
\(S(t)\) is fixed.

\subsection{Common nonlinear midpoint solve}

Both production propagators use the same SCEM fixed point. Given an input
midpoint density, build
\((S_{1/2},H_{\mathrm{eom},1/2},\omega_{1/2})\), construct a
trial propagator, obtain an output midpoint density, and iterate to
consistency. The convergence measure is
\begin{equation}
 r_P=
 \frac{
 \norm{S_{1/2}^{1/2}
 (P_{\mathrm{out}}-P_{\mathrm{in}})
 S_{1/2}^{1/2}}_F
 }{\max(N_e,1)}.
 \label{eq:scem-residual}
\end{equation}
The initial defaults are \(r_P\le10^{-10}\), at most 50 iterations, and
minimum damping 0.1. Start with a full Hermitian density update and reduce
mixing when the residual worsens. Record the iteration count, final density
residual, and an independent Hamiltonian residual. A failed midpoint is never
accepted; tolerances and \(\Delta t\) are never silently changed.

\subsection{Fixed-metric transport}

For bare LG/VG, \(S=S_0\) and \(\omega=0\). The midpoint generator is
\begin{equation}
 \mathcal A_{1/2}
 =-\frac{\ii}{\hbar}S^{-1}H_{1/2},
\end{equation}
which is \(S\)-skew-Hermitian. Its rational map can therefore act directly in
the nonorthogonal coordinates while preserving the \(S\) norm algebraically.
A L\"owdin transform is not required.

\subsection{Connection-aware site-parallel frame}

Write the analytically known diagonal site connection as \(\sigma(t)\).
Its exact parallel transport over one interval is
\begin{equation}
 G_\sigma(t_{n+1},t_n)
 =\exp\!\left[-\int_{t_n}^{t_{n+1}}\sigma(t)\,\dd t\right].
 \label{eq:site-transport}
\end{equation}
For the diagonal Abelian uniform-field case this is elementwise and depends
only on the known source. Pull the overlap, Hamiltonian, and residual
connection into this site-parallel frame. At the self-consistent midpoint,
the combined generator is
\begin{equation}
 \bar{\mathcal A}_{1/2}
 =\bar S_{1/2}^{-1}
 \left[-\bar\omega_{\mathrm{res},1/2}
       -\frac{\ii}{\hbar}\bar H_{1/2}\right].
 \label{eq:pulled-generator}
\end{equation}
The E1 contribution in Eq.~\eqref{eq:e1-split} is part of
\(\bar\omega_{\mathrm{res}}\), not an auxiliary length-gauge Hamiltonian.

\subsection{Rational exponential and metric correction}

The production default is the Pad\'e [2/2] map
\begin{equation}
 R_{22}(Z)=
 \left(I-\frac Z2+\frac{Z^2}{12}\right)^{-1}
 \left(I+\frac Z2+\frac{Z^2}{12}\right),
 \qquad Z=\Delta t\,\bar{\mathcal A}_{1/2}.
 \label{eq:pade22}
\end{equation}
Cayley [1/1] is retained as a diagnostic. These are systematically
improvable approximations to the frozen exponential\cite{Higham2008}, but the
complete midpoint method remains globally second order unless the
Hamiltonian/source quadrature and nonlinear construction are also raised in
order.

After transforming the raw interval map back to AO coordinates, impose the
cross-metric constraint by right Cholesky correction:
\begin{align}
 M&=U_{\mathrm{raw}}^\dagger S_{n+1}U_{\mathrm{raw}}
   =R_M^\dagger R_M,\nonumber\\
 S_n&=R_S^\dagger R_S,\nonumber\\
 Q&=R_M^{-1}R_S,\qquad
 U=U_{\mathrm{raw}}Q.
 \label{eq:cholesky-correction}
\end{align}
Then \(U^\dagger S_{n+1}U=S_n\) up to roundoff. Record the raw defect,
corrected defect, and \(\norm{Q-I}\). This is a constraint correction to an
already constructed physical transport; it is not propagation in an evolving
symmetrically orthogonalized L\"owdin basis.

Only roundoff-scale anti-Hermitian contamination may be symmetrized, under a
separately recorded threshold. A larger defect is a hard failure.

\section{Observable and storage model}

\subsection{Typed observable calculators}

An observable declares:

\begin{itemize}
\item the formulation capabilities and operator/source dependencies it needs;
\item whether it is evaluated at endpoints, midpoints, or events;
\item its independent integer-step schedule;
\item units, tensor shape, definition ID, and decomposition labels;
\item whether evaluation requires a Fock build or may reuse midpoint data.
\end{itemize}

Unsupported observables fail during simulation construction. Energy is not
computed every step by default. Energy, dipole/current, step diagnostics,
checkpoints, and dense-matrix snapshots have independent schedules, while
initial and final requested observables are always captured.

\subsection{Canonical artifacts}

The canonical numerical format is versioned HDF5:

\begin{center}
\begin{tabularx}{\textwidth}{@{}lX@{}}
\toprule
Artifact & Purpose\\
\midrule
\code{reference.h5} & Immutable backend-neutral prepared electronic reference\\
\code{trajectory.h5} & Canonical observable and diagnostic trajectory\\
\code{em\_source.h5} & Optional full-resolution/per-provider source history\\
\code{checkpoint\_<step>.h5} & Immutable restart state at an accepted boundary\\
\code{status.json} & Small atomic status record for detached monitoring\\
\bottomrule
\end{tabularx}
\end{center}

Top-level trajectory groups are
\begin{gather*}
 \texttt{/meta},\quad
 \texttt{/configuration},\quad
 \texttt{/reference},\quad
 \texttt{/time},\quad
 \texttt{/source},\\
 \texttt{/observables},\quad
 \texttt{/diagnostics},\quad
 \texttt{/events},\quad
 \texttt{/restart}.
\end{gather*}
Definition-specific observable paths are required; generic
\texttt{/current}, \texttt{/dipole}, or \texttt{/energy} datasets are
forbidden. Each stream owns step/time coordinates. Every dataset has unit and
dimension attributes. Scalar streams are uncompressed; dense/checkpoint arrays
are chunked on record boundaries with low-level gzip and shuffle.

Writers publish transactionally: write a \code{.partial} file, validate and
flush it, set an internal completion marker, and atomically rename it.
Completed artifacts are immutable. Controlled failures may publish
\code{trajectory.failed.h5} with \code{complete=false}; a hard crash leaves
only the partial artifact.

\subsection{Restart and lineage}

A checkpoint contains the dynamic state, occupations, global step/time,
accumulated work, relevant nonlinear continuation data, and applied-event IDs.
It contains no pickled PySCF or GPU object. Restart reconstructs the prepared
reference from normalized input and verifies all fingerprints.

A resumed segment first reconstructs, writes, and authenticates the repeated
boundary state \(n_r\) before advancing to \(n_r+1\). It records parent run
ID, checkpoint hash, and global step offset. Analysis deduplicates the shared
boundary only through this lineage. SIGINT/SIGTERM is handled at a safe
accepted-step boundary; a half-converged midpoint is never checkpointed as a
valid state.

\subsection{Semantic identity and provenance}

A deterministic simulation ID is the SHA-256 hash of losslessly normalized
scientific configuration, array dtypes/shapes/bytes, source and reference
fingerprints. Paths, timestamps, host names, and output schedules are
excluded. A unique run ID distinguishes execution attempts. Full software,
Git, dependency, CPU/GPU, and device provenance is stored. A requested GPU
run never falls back to CPU.

\section{Spectroscopy}

Casida linear response and real-time kick spectra are postprocessing
workflows. Casida roots, oscillator strengths, and transition directions are
stored as structured numerical data\cite{Casida1995}. A resonant comparison
selects a polarization-bright root by an explicit threshold or root choice and
recomputes it for every basis; otherwise an on-resonance/off-resonance
comparison would be physically ambiguous.

For a kick \(\vect\kappa\), the canonical molecular response is
\begin{equation}
 \alpha_{\alpha\beta}(\Omega)
 =\frac{\delta\widetilde\mu_\alpha(\Omega)}
        {\kappa_\beta}.
\end{equation}
The transform convention is
\begin{equation}
 \widetilde f(\Omega)
 =\int_0^T w(t)f(t)\ee^{+\ii\Omega t}\,\dd t.
 \label{eq:fft-convention}
\end{equation}
The full complex result, sign, normalization, time origin, quadrature, window,
damping, and zero padding are metadata, not hidden plotting choices. The
default window is rectangular. Optional exponential damping
\(\ee^{-\eta t}\) stores \(\eta\) in energy units. The intrinsic frequency
resolution is approximately \(2\pi/T\); zero padding interpolates the sampled
spectrum but does not improve this resolution. A current-domain transform
supplies an independent frequency-domain check.

\section{Reusable API and implementation boundaries}

\subsection{Static preparation and dynamic construction}

The top-level Python workflow is intentionally small:
\begin{verbatim}
reference = prepare_reference(reference_config)
reference.save("reference.h5")

simulation = build_simulation(
    reference=reference,
    source=physical_source,
    formulation=formulation_config,
    propagation=propagation_config,
    observables=observable_config,
    output=output_config,
)
result = run(simulation)
resumed = resume(checkpoint_path, output=child_output)
trajectory = load_trajectory(path)
\end{verbatim}

One prepared reference is reused by all formulations in a gauge comparison,
preventing repeated SCF work and guaranteeing a common initial problem.
References, configurations, operator bundles, sources, and formulation
definitions are immutable. Backend caches and preallocated scratch arrays
belong to a mutable \code{Workspace}; two simulations never share state or
workspace.

\subsection{Package domains}

\begin{longtable}{@{}p{0.34\textwidth}p{0.60\textwidth}@{}}
\toprule
Package & Responsibility\\
\midrule
\endhead
\code{aion.config} &
Immutable typed input, strict validation, normalization, units, and hashes.\\
\code{aion.backends} &
Array protocol, CPU/GPU implementations, device residency, and workspaces.\\
\code{aion.electronic\_structure} &
PySCF RKS preparation, grids, reference reconstruction, core and derived
operator bundles.\\
\code{aion.electromagnetism} &
Physical sources, gauge representations, node/link samples, envelopes,
compiled time series, and events.\\
\code{aion.formulations} &
Bare LG/VG and P0/P0+E1 implementations of
\((S,\hH,\omega)\), formulation observables, and support checks.\\
\code{aion.propagation} &
State protocols, common nonlinear SCEM, fixed-metric and connection-aware
transport.\\
\code{aion.observables} &
Typed dipole, current, mechanical-current, energy, power, continuity, and Ward
calculators.\\
\code{aion.io} &
HDF5 schemas, transactional writers, checkpoints, lineage, readers, and
export.\\
\code{aion.workflows} &
Preparation, simulation building, run/resume, comparisons, Casida, and kick
spectra.\\
\bottomrule
\end{longtable}

Dependencies are one-way: formulations consume references and sources;
propagators consume the formulation protocol; observables inspect states and
formulations; I/O receives records but never enters the physics. The CLI
\code{aion prepare|run|resume|inspect|export} is a thin caller of the same
typed API.

\subsection{Prepared reference content}

The immutable \code{reference.h5} contains normalized molecule/basis/
functional/grid input, converged coefficients and occupations, ground-state
density, nuclear data, EM origin, exact quadrature-grid coordinates and
weights, and the core operator bundle
\[
  \{S,T,V_{\mathrm{nuc}},\vect r,\vect p,
    \text{nuclear data},f\}.
\]
It does not store AO values on the grid and never pickles live objects.
Loading reconstructs PySCF/backend models and verifies fingerprints.
Atomic projectors/topology and P0/E1 tensors are independently versioned
derived bundles with their own invalidation boundaries.

\section{Present implementation versus the 0.2 target}

\begin{statusbox}{Current-tree assessment after WP1}
The research draft was archived and verified on 8 September 2026. Its
molecule-specific campaigns, runners, postprocessing, historical notes, and
large results were moved with checksum manifests to the separate
\path{CALCULATIONS} repository. WP1 then removed the active flat prototype
modules and established the clean 0.2 package contracts. Draft equations and
kernels remain available through the immutable archive and calculation
repository, but none is silently reachable through the 0.2 API.
\end{statusbox}

The reusable layers implemented by WP1 are:

\begin{longtable}{@{}p{0.31\textwidth}p{0.63\textwidth}@{}}
\toprule
Current package & Implemented contract\\
\midrule
\endhead
\code{aion.config} &
Frozen reference and simulation dataclasses, atomic-unit origin and fixed-grid
types, discriminated formulations/sources/events, strict TOML resolution, and
lossless canonical SHA-256 identities.\\
\code{aion.io} &
Independent semantic schema versions, structural validators and skeletons for
reference, trajectory, source-history, and checkpoint HDF5 artifacts, plus the
strict non-authoritative \code{status.json} record. Transactional writers are
not yet implemented.\\
\code{aion.observables} &
Definition-specific immutable observable metadata and sampled records. The
physical dipole, current, energy, power, Ward, and continuity calculators are
not yet implemented.\\
\code{aion.workflows} and \code{aion} &
Typed top-level interfaces and CLI commands. Numerical entry points fail
explicitly with their owning future work package; configuration validation and
HDF5 header inspection are functional.\\
\bottomrule
\end{longtable}

Important missing reusable layers are:

\begin{itemize}
\item a unified immutable prepared reference and compiled-source/simulation
  builder behind the established configuration contracts;
\item formulation-owned physical primary and diagnostic observable
  calculators behind the established record types;
\item the generic mechanical-current operator API;
\item the complete qualified energy ledger and independent analytic
  \(\dd U/\dd t\);
\item a single shared SCEM nonlinear engine behind both transports;
\item transactional HDF5 writers with exact restart lineage;
\item completed numerical CLI commands and reusable spectroscopy outputs.
\end{itemize}

At present, some comparison scripts reconstruct missing VG linear and
diamagnetic kinetic energies, calculate bare-LG momentum current, or assemble
covariant mechanical diagnostics after propagation. Plotting scripts also
encode observable definitions. In 0.2 these calculations belong to
\code{aion.observables}; campaigns select and plot stored definitions but do
not rederive them.

\section{Validation requirements}

\subsection{Hierarchy}

The validation fixtures have distinct purposes:

\begin{description}[leftmargin=32mm,style=nextline]
\item[Synthetic matrices]
Exact signs, Hermiticity, metric compatibility, graph incidence/orientation,
gauge covariance, rational-map algebra, and cross-metric transport.
\item[H\(_2\)]
Zero field, molecular symmetry, bare-gauge smoke tests, exact kicks, and
compact timestep convergence.
\item[LiH]
Nonzero E1 response, gauge/origin transformations, variational current,
continuity, Ward energy, connection transport, restart, and CPU/GPU parity.
\item[NH\(_3\)]
Integrated release validation only: aug-cc-pVTZ, five driven cycles plus five
field-free cycles, all four formulations, accepted step and a shorter
half-step interval.
\end{description}

CO, ammonia/PNA production studies, and basis-convergence campaigns are
external scientific calculations, not unit-test decisions. Existing draft
trajectories are evidence, not golden truth.

\subsection{Independent internal checks}

The principal assertions are:

\begin{enumerate}
\item \(P=P^\dagger\), electron number and occupied-state metric
  normalization;
\item Eq.~\eqref{eq:metric-compatibility} and cross-metric preservation;
\item explicit gauge-function transformations and EMF invariance;
\item translation of the EM origin, including electronic/nuclear
  cancellation for neutral systems;
\item Eq.~\eqref{eq:cov-continuity} on every node and consistency of pair
  signs with Eq.~\eqref{eq:graph-current};
\item directional derivatives of the variational source action versus
  Eq.~\eqref{eq:source-current};
\item analytic Eq.~\eqref{eq:j-mu}, independent of finite differences;
\item energy-component sums and analytic Eq.~\eqref{eq:analytic-energy-rate}
  versus source power;
\item approximately second-order reduction of trajectory and Ward errors
  under step halving;
\item uninterrupted versus restarted trajectory identity;
\item CPU/GPU parity for state, observables, energy components, invariants,
  and backend residency.
\end{enumerate}

Full-array pinned trajectories are not the primary oracle. Selected numerical
endpoints are retained only when they provide an independent protection.

\subsection{GPU and performance policy}

Every milestone that changes executable numerical physics must run its
applicable tests on the physical GPU. The execution sandbox may hide the
device; this is handled by requesting elevated execution, never by omitting
GPU tests or substituting CPU-only tests. A requested GPU backend fails if
unavailable.

The fast test suite uses up to eight workers with one BLAS thread each.
Molecular CPU tests run serially with up to eight BLAS threads, and GPU tests
run serially. Preparation, Fock construction, steady step time, memory, I/O,
and host transfers are measured separately. An unexplained steady-step
regression of roughly ten percent or more is flagged, but correctness and
physics gates take priority.

\section{Conclusions}

The bare and covariant formulations solve related physical problems but do
not share interchangeable current or energy definitions. Bare LG uses the
analytic dipole derivative as its primary current; bare VG uses paramagnetic
plus diamagnetic mechanical current and requires the complete mechanical
kinetic energy; P0+E1 uses the variational source current implied by its
projected connection. A common ambient-projected mechanical current is an
important diagnostic, not a replacement for the covariant source current.

The difficult numerical feature of P0+E1 is not ignorance of the connection:
the connection is analytically known from prescribed fields and static matrix
elements. The issue is faithful finite-step transport between distinct
metrics while the nonlinear Hamiltonian is determined at the midpoint. Exact
site parallel transport, a combined residual generator, Pad\'e midpoint
propagation, and a right Cholesky cross-metric correction address that geometry
without an evolving L\"owdin frame.

The 0.2 architecture makes these physics definitions reusable and
auditable. It separates immutable reference preparation, physical sources,
formulations, propagation, observables, and storage; it records qualified
energy/current definitions; and it treats conservation, Ward identities,
convergence, restart, and CPU/GPU parity as release gates. The preservation and
repository-separation gate WP0 and the non-numerical contract gate WP1 are
complete. Numerical implementation begins with WP2 and remains subject to its
separate acceptance gates.

\section*{Rebuilding this document}

From the \path{aion/docs} directory, the PDF is rebuilt and all references are
resolved with the single shell command
\begin{verbatim}
for pass in 1 2 3; do \
  /usr/bin/pdflatex -interaction=nonstopmode \
  -halt-on-error theory_and_implementation.tex; \
done
\end{verbatim}
The generated auxiliary, log, outline, and table-of-contents files are build
artifacts; the tracked deliverables are the TeX source and PDF.

\begin{thebibliography}{99}

\bibitem{RungeGross1984}
E.~Runge and E.~K.~U. Gross,
\emph{Density-Functional Theory for Time-Dependent Systems},
\emph{Physical Review Letters} \textbf{52}, 997--1000 (1984),
\href{https://doi.org/10.1103/PhysRevLett.52.997}
{doi:10.1103/PhysRevLett.52.997}.

\bibitem{Casida1995}
M.~E. Casida,
\emph{Time-dependent density functional response theory for molecules},
in \emph{Recent Advances in Density Functional Methods, Part I},
edited by D.~P. Chong (World Scientific, 1995), pp.~155--192.

\bibitem{Pemmaraju2018}
C.~D. Pemmaraju, F.~D. Vila, J.~J. Kas, S.~A. Sato, J.~J. Rehr, K.~Yabana,
and D.~Prendergast,
\emph{Velocity-gauge real-time TDDFT within a numerical atomic orbital basis
set},
\emph{Computer Physics Communications} \textbf{226}, 30--38 (2018),
\href{https://doi.org/10.1016/j.cpc.2018.01.013}
{doi:10.1016/j.cpc.2018.01.013}.

\bibitem{ArtachoORegan2017}
E.~Artacho and D.~D. O'Regan,
\emph{Quantum mechanics in an evolving Hilbert space},
\emph{Physical Review B} \textbf{95}, 115155 (2017),
\href{https://doi.org/10.1103/PhysRevB.95.115155}
{doi:10.1103/PhysRevB.95.115155}.

\bibitem{PBE1996}
J.~P. Perdew, K.~Burke, and M.~Ernzerhof,
\emph{Generalized Gradient Approximation Made Simple},
\emph{Physical Review Letters} \textbf{77}, 3865--3868 (1996),
\href{https://doi.org/10.1103/PhysRevLett.77.3865}
{doi:10.1103/PhysRevLett.77.3865}.

\bibitem{Sun2018}
Q.~Sun, T.~C. Berkelbach, N.~S. Blunt, G.~H. Booth, S.~Guo, Z.~Li,
J.~Liu, J.~D. McClain, E.~R. Sayfutyarova, S.~Sharma, S.~Wouters, and
G.~K.-L. Chan,
\emph{PySCF: the Python-based simulations of chemistry framework},
\emph{WIREs Computational Molecular Science} \textbf{8}, e1340 (2018),
\href{https://doi.org/10.1002/wcms.1340}
{doi:10.1002/wcms.1340}.

\bibitem{Higham2008}
N.~J. Higham,
\emph{Functions of Matrices: Theory and Computation}
(Society for Industrial and Applied Mathematics, 2008),
\href{https://doi.org/10.1137/1.9780898717778}
{doi:10.1137/1.9780898717778}.

\end{thebibliography}

\end{document}
